\chapter{Si venías de Windows}\label{cap:c-equivalencias}

\begin{center}
\begin{tabularx}{\textwidth}{XX}
\toprule
En Windows & En Lázaro \\
\midrule
Menú Inicio & Menú \menu{Aplicaciones} (arriba a la izquierda en XFCE) \\
Explorador de archivos & Gestor de archivos (Thunar en XFCE) \\
Panel de control / Configuración & \menu{Aplicaciones > Configuración >
  Gestor de configuración} \\
Microsoft Store & La tienda de aplicaciones (Software o Discover) \\
Windows Update & El aviso de actualizaciones, \menu{Actualizar el sistema} \\
Administrador de tareas & Btop (\texttt{btop} en la terminal) o el monitor
  del sistema de tu escritorio \\
Word, Excel, PowerPoint & LibreOffice Writer, Calc, Impress (u OnlyOffice) \\
Edge, Chrome & Firefox (Chrome y Brave, opcionales) \\
Adobe Reader & Visor de documentos (Evince) \\
Reproductor de Windows Media & mpv o VLC \\
Fotos & Ristretto (visor de fotos de XFCE) \\
Bloc de notas & Mousepad (editor de texto de XFCE) \\
Recortes / Imprimir pantalla & Captura de pantalla (\keys{ImprPant}) \\
Calculadora & Galculator \\
WinRAR, 7-Zip & Clic derecho, \menu{Extraer aquí} \\
\enquote{Ejecutar como administrador} & \texttt{doas} delante de la orden \\
Símbolo del sistema, PowerShell & La terminal \\
Antivirus & No hace falta: los programas vienen de fuentes firmadas y se
  actualizan solos. Mantén el sistema al día. \\
Unidad C: & La carpeta personal, para tus archivos \\
\bottomrule
\end{tabularx}
\end{center}
